\section*{\Huge Acknowledgement}
\addcontentsline{toc}{chapter}{Acknowledgement}

This specialized project is the first time we carried a piece of work from an open research question to a working system, and we did not do it alone.

First and foremost, we express our deepest gratitude to our supervisors, Assoc. Prof. Quan Thanh Tho, Ph.D and Assoc. Prof. Nguyen Duc Anh, Ph.D. They encouraged us to explore a young and fast-moving field, helped us turn a broad idea into precise research questions, and reminded us that every design choice needs a reason and every claim needs evidence. Their patient feedback on our ideas, our design and our writing shaped both this report and the way we approach problems.

We are grateful to the lecturers of the Faculty of Computer Science and Engineering, Ho Chi Minh City University of Technology, Vietnam National University Ho Chi Minh City. The knowledge of software engineering, requirements engineering and artificial intelligence that we gained in their courses formed the foundation of this project. We also thank the researchers and the open-source communities whose papers, models and tools we built upon; their willingness to share their work allowed a small student team to stand on solid ground.

To our friends and classmates, thank you for the discussions that sharpened our ideas, for trying our prototype and telling us honestly what did not work, and for your encouragement during the long evenings of debugging and writing. We thank each other as well: the project taught us how to share work, review each other's contributions, disagree constructively and keep going together when progress was slow.

Finally, we owe our greatest debt to our families. Their unconditional love, patience and quiet sacrifices have supported us throughout our years of study, and they believed in us even when we doubted ourselves. This work is dedicated to them.

Despite our efforts, this report surely still has shortcomings, and we sincerely welcome the comments of our supervisors and the committee to improve it in the capstone (thesis) phase.
